\section{Відповіді на контрольні запитання}

\begin{enumerate}
    \item \textbf{Що таке успадкування в об'єктно-орієнтованому програмуванні?}
    Це фундаментальний принцип ООП, який дозволяє створювати новий (дочірній) клас 
    на основі вже існуючого (батьківського), автоматично переймаючи його захищені 
    та публічні поля і методи для повторного використання коду та розширення 
    функціональності.

    \item \textbf{Яке ключове слово використовується для успадкування класів у Java?}
    Для реалізації успадкування класів у мові програмування Java використовується 
    ключове слово \texttt{extends}.

    \item \textbf{Чим відрізняється базовий клас від дочірнього?}
    Базовий (батьківський) клас описує загальні універсальні характеристики сутностей, 
    тоді як дочірній (похідний) клас конкретизує поведінку, додає власні унікальні 
    поля/методи та перевизначає логіку суперкласу.

    \item \textbf{Чи може клас у Java успадковувати декілька класів одночасно?}
    Ні, в Java реалізовано лише поодиноке успадкування для класів (конструкція з 
    кількома батьківськими класами не підтримується), проте один клас може реалізовувати 
    декілька інтерфейсів одночасно.

    \item \textbf{Для чого використовується ключове слово super?}
    Ключове слово \texttt{super} використовується для звернення до елементів батьківського
     класу: для виклику його конструктора на початку конструктора дочірнього класу 
     (наприклад, \texttt{super(brand, model)}) або для виклику оригінальних методів 
     суперкласу.

    \item \textbf{Що таке перевизначення методу?}
    Це операція, за якої дочірній клас змінює внутрішню реалізацію (тіло методу) 
    успадкованого від батьківського класу методу при повному збереженні його імені, 
    аргументів та типу повертаємого значення.

    \item \textbf{Для чого використовується анотація @Override?}
    Вона рекомендується для позначення перевизначеного методу та вказує компілятору Java 
    на перевірку відповідності сигнатур. Це захищає від помилок друку — якщо сигнатури 
    методів базового та дочірнього класів не збігаються, компілятор видасть помилку.

    \item \textbf{Що таке абстрактний клас?}
    Це клас, оголошений із ключовим словом \texttt{abstract}, який призначений для 
    створення спільної основи для інших класів. Він може містити як звичайні методи 
    з реалізацією, так і абстрактні методи без тіла.

    \item \textbf{Чи можна створити безпосередньо об'єкт абстрактного класу?}
    Ні, створити безпосередньо об'єкт абстрактного класу неможливо (наприклад, виклик 
    \texttt{new ComputerDevice()} призведе до помилки компіляції).

    \item \textbf{Що таке абстрактний метод?}
    Це метод, який оголошується без тіла (без конкретної реалізації, завершується 
    крапкою з комою), задаючи лише сигнатуру, яку всі конкретні дочірні класи обов'язково 
    повинні реалізувати.

    \item \textbf{Які елементи може містити абстрактний клас?}
    Абстрактний клас може містити поля (змінні стану), конструктори, звичайні методи з 
    реалізацією, а також абстрактні методи без реалізації.

    \item \textbf{Що таке інтерфейс у Java?}
    Це повністю абстрактна конструкція, створена за допомогою ключового слова 
    \texttt{interface}, яка визначає певний набір методів (контракт поведінки), що повинні 
    бути реалізовані класом.

    \item \textbf{Яке ключове слово використовується для реалізації інтерфейсу?}
    Клас реалізує інтерфейс за допомогою ключового слова \texttt{implements}.

    \item \textbf{Чи може один клас реалізовувати декілька інтерфейсів?}
    Так, однією з важливих переваг інтерфейсів є можливість реалізувати декілька інтерфейсів 
    одним конкретним класом через кому.

    \item \textbf{У чому полягає основна відмінність між абстрактним класом та інтерфейсом?}
    Абстрактний клас доцільно використовувати тоді, коли класи мають спільний стан, поля та 
    частково спільну реалізацію (зв'язок сутностей). Інтерфейс доцільно використовувати для 
    визначення певної можливості або набору операцій, які можуть бути реалізовані різними, 
    навіть не пов'язаними між собою класами.

    \item \textbf{Що таке поліморфізм?}
    Це принцип ООП, який дає змогу працювати з об'єктами різних дочірніх класів через 
    спільний базовий тип (суперклас або інтерфейс), забезпечуючи виконання потрібної 
    реалізації методу на етапі виконання програми.

    \item \textbf{Для чого можна використовувати інтерфейс як тип змінної?}
    Це дозволяє створювати універсальний, гнучкий код, який абстрагується від конкретних 
    класів у кучі (Heap) та працює з усіма об'єктами, що реалізують відповідний інтерфейсний 
    контракт (наприклад, \texttt{Connectable officeLaptop = new Laptop(...)}).

    \item \textbf{Що відбувається під час виклику перевизначеного методу через змінну базового типу?}
    Завдяки механізму динамічного поліморфізму, під час виконання програми віртуальна машина 
    Java (JVM) аналізує фактичний тип об'єкта, на який посилається змінна, і викликає саме 
    перевизначену реалізацію дочірнього класу.

    \item \textbf{Для чого в об'єктно-орієнтованому програмуванні використовують абстракцію?}
    Для виділення головних, спільних та визначальних характеристик об'єктів певної предметної 
    області й відсікання другорядних деталей реалізації, що допомагає проектувати архітектуру 
    складних систем.

    \item \textbf{Які переваги дає використання успадкування та інтерфейсів під час розроблення програм?}
    Успадкування дозволяє уникати дублювання коду та будувати чіткі ієрархії пов'язаних класів, 
    а інтерфейси дозволяють визначати спільні гнучкі правила взаємодії між різними класами, 
    забезпечуючи легкість розширення та масштабування програми.
\end{enumerate}